\scanpage{11}
\mspage{10}

\uline{Conditions particulières sur les lieux.}

\uline{NB} Pour qu'un espace compact triangulable admette une structure de ficelle, il f. et s. qu'il soit connexe, et non réduit à un point. C'est gros! Il y a alors $=$ une structure de ficelle réduite. [NB. Si $X$ est un segment cette str. est unique. C'est le seul cas d'exception!]

\margin{\uline{NB} les seules \deleted{structures}\unsure{de|\deleted{de}}\inserted{ficelles} sur $[0,1]$ sont données par $\{A,B\}$ où $A=[0,\alpha]$, $B=[\beta,1]$ avec $0\le\alpha<\unsure{\beta'|\beta}\le1$. \deleted{\unsure{NB}}}

On peut considérer relation d'équivalence entre structures de ficelles \deleted{sur $X$} (resp. ficelles orientées) sur $X$: \deleted{\unsure{p}} celle engendrée par la relation d'ordre (cas orienté)
\deleted{ou \unsure{$(A\subset A'\ \text{ou}\ A\subset B')$ et $(B\subset A'\ \unsure{,|\text{ou}}\ B\subset B')$}}
\[
A\subset A'\quad\text{et}\quad B\subset B'.
\]

\margin{Il est \unsure{raisonnable} de supposer les extrémités $A,B$ (\unsure{de même} $A',B'$) connexes voire, \unsure{contractiles}.}
\deleted{Oui}
\deleted{\uline{NB} Si $(X,\{A,B\})$ est une ficelle, et si $A',B'$ sont des fermés régulièrement immergés tels que $A\subset A'$, $B\subset B'$, $A'\cap B'=\varnothing$, alors $(X,\{A',B'\})$ est une ficelle.}
\margin{Faux $\longrightarrow$ idiot!}
\margin{(\unsure{Connexité})}

\deleted{donc}\inserted{soupçonne} que deux structures de ficelle orientée sur un $X$ triangulable sont toujours équivalentes.
\margin{Faux déjà dans le cas du segment!}

\rotatebox{90}{\reflectbox{\unsure{EXTRA STRONG}}}
\rotatebox{90}{\reflectbox{\unsure{EXTRA STRONG}}}
\mspage{10}
